\chapter{The Statistical Ising Model and Universality}
\label{ch:volume-ii-ising-universality}

\section{The Finite Lattice Ensemble}

Fix a spatial dimension \(D\geq 1\).  A finite Ising system is specified by a
finite subset \(\Lambda\subset \mathbb Z^D\), a set \(E(\Lambda)\) of
nearest-neighbor unordered pairs, and spin variables
\[
  s_x\in\{-1,+1\},
  \qquad x\in\Lambda .
\]
For ferromagnetic coupling \(J>0\), external magnetic field \(h\), and inverse
temperature \(\beta=T^{-1}\) in units with Boltzmann's constant equal to one,
the energy of a spin configuration \(s=\{s_x\}_{x\in\Lambda}\) is
\[
  H_\Lambda(s)
  =
  -J\sum_{\langle x y\rangle\in E(\Lambda)}s_xs_y
  -h\sum_{x\in\Lambda}s_x .
\]
The finite-volume partition function and thermal expectation of an observable
\(F:\{-1,+1\}^{\Lambda}\to\mathbb C\) are
\[
  Z_\Lambda(\beta,h)
  =
  \sum_{s_x=\pm1}
  \exp\!\left[-\beta H_\Lambda(s)\right],
  \qquad
  \vev{F}_{\Lambda,\beta,h}
  =
  Z_\Lambda(\beta,h)^{-1}
  \sum_{s_x=\pm1}F(s)\,\ee^{-\beta H_\Lambda(s)} .
\]
At \(h=0\) the model has the global \(\mathbb Z_2\) symmetry
\(s_x\mapsto -s_x\) for all \(x\).

The same finite probability space can be written in diagonal operator notation.
Let
\[
  V_x=\operatorname{span}\{\ket{+}_x,\ket{-}_x\},
  \qquad
  \Hilb_\Lambda=\bigotimes_{x\in\Lambda}V_x ,
\]
and define commuting spin operators by
\(\hat S_x\ket{s}_\Lambda=s_x\ket{s}_\Lambda\), where
\(\ket{s}_\Lambda=\bigotimes_x\ket{s_x}_x\).  Then
\[
  \hat H_\Lambda
  =
  -J\sum_{\langle x y\rangle}\hat S_x\hat S_y
  -h\sum_x\hat S_x,
  \qquad
  Z_\Lambda(\beta,h)
  =
  \operatorname{Tr}_{\Hilb_\Lambda}\ee^{-\beta\hat H_\Lambda}.
\]
The space \(\Hilb_\Lambda\) in this display is a finite-dimensional
representation of a classical statistical ensemble.  A continuum Hilbert space,
when constructed from the scaling limit of Euclidean correlators, is a further
object obtained after the limit.

\begin{figure}[H]
  \centering
  \resizebox{0.98\textwidth}{!}{%
  \begin{tikzpicture}[x=1cm,y=1cm,every node/.style={font=\scriptsize}]
    \begin{scope}
      \node[font=\small] at (1.5,3.2) {classical configurations};
      \foreach \x in {0,1,2,3} {
        \foreach \y in {0,1,2} {
          \fill[black] (\x,\y) circle (1.5pt);
        }
      }
      \foreach \x in {0,1,2} {
        \foreach \y in {0,1,2} {
          \draw[gray!60] (\x,\y)--(\x+1,\y);
        }
      }
      \foreach \x in {0,1,2,3} {
        \foreach \y in {0,1} {
          \draw[gray!60] (\x,\y)--(\x,\y+1);
        }
      }
      \foreach \x/\y/\spin in {0/0/+,1/0/-,2/0/-,3/0/+,0/1/-,1/1/+,2/1/-,3/1/-,0/2/+,1/2/+,2/2/-,3/2/+} {
        \node[blue!70!black] at (\x,\y+0.28) {\(\spin\)};
      }
      \node[align=center] at (1.5,-0.65)
        {\(s_x=\pm1\), \(x\in\Lambda\subset\mathbb Z^D\)};
    \end{scope}

    \begin{scope}[shift={(5.0,0)}]
      \node[font=\small] at (0,3.2) {Boltzmann weight};
      \node[draw=black, rounded corners, align=center, inner sep=7pt] at (0,1.55)
        {\(\displaystyle
        \mathbb P_{\beta,h}(s)
        =
        Z_\Lambda^{-1}\ee^{-\beta H_\Lambda(s)}
        \)\\[2pt]
        \(\displaystyle
        H_\Lambda=-J\sum_{\langle x y\rangle}s_xs_y-h\sum_xs_x
        \)};
    \end{scope}

    \begin{scope}[shift={(9.5,0)}]
      \node[font=\small] at (0,3.2) {diagonal trace notation};
      \node[draw=black, rounded corners, align=center, inner sep=7pt] at (0,1.55)
        {\(\displaystyle
        \Hilb_\Lambda=\bigotimes_{x\in\Lambda}V_x
        \)\\[2pt]
        \(\displaystyle
        Z_\Lambda=\operatorname{Tr}_{\Hilb_\Lambda}
        \ee^{-\beta\hat H_\Lambda}
        \)};
      \node[align=center, blue!70!black] at (0,-0.65)
        {all \(\hat S_x\) commute};
    \end{scope}

    \draw[-{Latex[length=2mm]}, thick] (3.55,2.62)--(4.35,2.62);
    \draw[-{Latex[length=2mm]}, thick] (6.9,2.62)--(7.65,2.62);
  \end{tikzpicture}
  }
  \caption{The statistical Ising model is a probability measure on lattice spin
  configurations.  Diagonal operators provide an exact finite-volume trace
  notation for the same measure.}
  \label{fig:volume-ii-ising-ensemble}
\end{figure}

For a finite set of energy eigenstates \(n\) with energies \(E_n\), the
canonical weights
\[
  p_n=\frac{1}{Z}\ee^{-\beta E_n},
  \qquad
  Z=\sum_n\ee^{-\beta E_n},
\]
are obtained by extremizing the Shannon entropy
\[
  S[p]=-\sum_n p_n\log p_n
\]
subject to \(\sum_np_n=1\) and \(\sum_np_nE_n=U\).  Indeed, the stationary
point of
\[
  S[p]+\alpha\Bigl(\sum_np_n-1\Bigr)
  +\beta\Bigl(U-\sum_np_nE_n\Bigr)
\]
has \(\log p_n=\alpha-1-\beta E_n\), hence the displayed Boltzmann form.

The entropy also has a simple monotonicity statement for symmetric probability
evolution.  If a finite distribution \(p_n(t)\) obeys
\[
  \dot p_n=\sum_m v_{nm}(p_m-p_n),
  \qquad
  v_{nm}=v_{mn}\geq0 ,
\]
then \(\sum_n\dot p_n=0\) and
\[
  \frac{\dd S}{\dd t}
  =
  \frac12\sum_{m,n}v_{mn}(p_m-p_n)(\log p_m-\log p_n)\geq0 .
\]
This is the finite-state entropy-increase mechanism underlying thermal
equilibration in the idealized symmetric-rate case.

\section{Spin Correlators and Phases}

The basic observables are spin correlation functions.  At \(h=0\), for
translation-invariant infinite-volume limits, write
\[
  C_\beta(x_1,\ldots,x_n)
  =
  \vev{s_{x_1}\cdots s_{x_n}}_\beta .
\]
The two-point function at large separation is often summarized by
\[
  f_\beta(r)
  \simeq
  \vev{s_0s_x}_\beta,
  \qquad r=|x|.
\]
For \(D=2,3\) and ferromagnetic nearest-neighbor coupling, there is a critical
temperature \(T_c\).  In the ordered region \(T<T_c\), pure infinite-volume
states selected by boundary conditions or by taking \(h\to0^\pm\) after the
thermodynamic limit have nonzero magnetization
\[
  m_\pm(T)
  =
  \lim_{h\to0^\pm}\lim_{\Lambda\nearrow\mathbb Z^D}
  \vev{s_0}_{\Lambda,\beta,h},
  \qquad m_+(T)=-m_-(T)\neq0 .
\]
The full two-point function tends to \(m_\pm(T)^2\) at large distance in such a
phase.  The connected correlator
\[
  G_\beta(x)
  =
  \vev{s_0s_x}_\beta-\vev{s_0}_\beta\vev{s_x}_\beta
\]
decays at large \(|x|\) away from criticality.  In the disordered region
\(T>T_c\),
\[
  G_\beta(x)\sim A(T)\,\ee^{-|x|/\xi(T)}
\]
up to powers of \(|x|\), where \(\xi(T)\) is the correlation length.  One
precise definition, when the limit exists along a coordinate direction, is
\[
  \xi(T)^{-1}
  =
  -\lim_{r\to\infty}\frac1r\log |G_\beta(r e_1)|.
\]
At \(T=T_c\), the correlation length diverges and the spin two-point function
has power-law asymptotics
\[
  \vev{s_0s_x}_{\beta_c}
  \sim
  \frac{A_\sigma}{|x|^{2\Delta_\sigma}} .
\]
The exponent \(\Delta_\sigma\) is the scaling dimension of the continuum spin
field.

\begin{figure}[H]
  \centering
  \resizebox{0.92\textwidth}{!}{%
  \begin{tikzpicture}[x=1cm,y=1cm,every node/.style={font=\scriptsize}]
    \draw[-{Latex[length=2mm]}, thick] (0,0)--(7.2,0) node[below] {\(r\)};
    \draw[-{Latex[length=2mm]}, thick] (0,0)--(0,3.1) node[left] {\(f_\beta(r)\)};
    \draw[green!50!black, very thick]
      (0,2.45)..controls (1.4,2.15) and (2.7,1.95)..(4.1,1.92)
      ..controls (5.0,1.90) and (6.0,1.95)..(6.7,1.96);
    \draw[blue!75!black, very thick]
      (0,2.45)..controls (1.2,1.7) and (2.4,1.10)..(3.8,0.72)
      ..controls (5.0,0.43) and (6.0,0.28)..(6.7,0.20);
    \draw[purple!75!black, very thick]
      (0,2.45)..controls (0.5,1.45) and (0.95,0.55)..(1.75,0.22)
      ..controls (2.5,0.06) and (3.4,0.03)..(4.8,0.02);
    \node[green!50!black] at (5.2,2.25) {\(T<T_c\)};
    \node[blue!75!black] at (5.1,0.75) {\(T=T_c\)};
    \node[purple!75!black] at (2.3,0.62) {\(T>T_c\)};
    \node[green!50!black, align=center] at (2.2,2.35)
      {ordered\\plateau};
    \node[blue!75!black, align=center] at (5.7,1.28)
      {\(f_\beta(r)\sim r^{-2\Delta_\sigma}\)};
    \node[purple!75!black, align=center] at (4.6,0.34)
      {\(\exp[-r/\xi(T)]\)};
  \end{tikzpicture}
  }
  \caption{Long-distance spin correlations distinguish the ordered,
  critical, and disordered regimes.  In the ordered phase the full two-point
  function approaches the square of the magnetization; the connected correlator
  still decays away from the critical point.}
  \label{fig:volume-ii-ising-correlator-regimes}
\end{figure}

The values relevant for the first Ising universality classes are
\[
\begin{array}{c|cc}
  D & \Delta_\sigma & \nu \\ \hline
  2 & \frac18 & 1\\[2pt]
  3 & 0.518\ldots & 0.62998\ldots
\end{array}
\]
where \(\nu\) is defined by the divergence
\[
  \xi(T)\sim \xi_0\,|T-T_c|^{-\nu}
  \qquad (T\to T_c).
\]

\section{The Scaling Limit}

The continuum Ising field theory is defined from large-distance lattice
correlators near criticality.  Let \(L\) be a large dimensionless length in
lattice units, and let \([Lx]\in\Lambda\) denote a nearest lattice point to
\(Lx\), where \(x\in\mathbb R^D\) is fixed.  Choose a sequence of temperatures
\(T(L)\to T_c\) such that
\[
  \frac{L}{\xi(T(L))}\longrightarrow m\geq0 .
\]
The number \(m\) is the mass gap of the continuum scaling limit in the units in
which the physical position is \(x\).  A spin-field renormalization factor
\(Z_\sigma(L,T)>0\) is chosen so that the limit
\[
  \lim_{\substack{T\to T_c,\ L\to\infty\\ L/\xi(T)\to m}}
  Z_\sigma(L,T)^{-n/2}
  \vev{s_{[Lx_1]}\cdots s_{[Lx_n]}}_{\beta(T)}
  =
  \vev{\sigma(x_1)\cdots\sigma(x_n)}_{\operatorname{Ising},m}
\]
exists for separated points.  At \(m=0\), equivalently at
\(\xi(T)/L\to\infty\), the limiting theory is the Ising conformal field
theory.  At criticality \(Z_\sigma(L,T_c)\) may be chosen proportional to
 \(L^{-2\Delta_\sigma}\), so the rescaled spin insertions have finite
correlators.

\begin{figure}[H]
  \centering
  \resizebox{0.96\textwidth}{!}{%
  \begin{tikzpicture}[x=1cm,y=1cm,every node/.style={font=\scriptsize}]
    \begin{scope}
      \node[font=\small] at (1.7,3.05) {lattice insertions};
      \draw[step=0.55, gray!45] (0,0) grid (3.3,2.2);
      \foreach \x/\y in {0.55/1.65,1.65/1.1,2.75/0.55} {
        \fill[blue!70!black] (\x,\y) circle (2.4pt);
      }
      \node[align=center] at (1.65,-0.45)
        {\([Lx_i]\) nearest lattice points};
    \end{scope}

    \draw[-{Latex[length=2mm]}, thick] (3.9,1.1)--(5.1,1.1)
      node[midway, above, align=center] {\(L\to\infty\)\\\(T\to T_c\)};

    \begin{scope}[shift={(5.8,0)}]
      \node[font=\small] at (1.65,3.05) {continuum fields};
      \draw[-{Latex[length=2mm]}, gray!70] (0,0)--(3.35,0) node[below] {\(x^1\)};
      \draw[-{Latex[length=2mm]}, gray!70] (0,0)--(0,2.35) node[left] {\(x^2\)};
      \foreach \x/\y in {0.55/1.65,1.65/1.1,2.75/0.55} {
        \fill[purple!75!black] (\x,\y) circle (2.4pt);
        \node[purple!75!black, above right=-1pt] at (\x,\y) {\(\sigma(x_i)\)};
      }
      \node[align=center, blue!70!black] at (1.65,-0.45)
        {\(L/\xi(T)\to m\)};
    \end{scope}

    \node[draw=black, rounded corners, align=center, inner sep=7pt]
      at (5.0,-1.45)
      {\(\displaystyle
        Z_\sigma^{-n/2}\,
        \vev{s_{[Lx_1]}\cdots s_{[Lx_n]}}
        \longrightarrow
        \vev{\sigma(x_1)\cdots\sigma(x_n)}_{\operatorname{Ising},m}
      \)};
  \end{tikzpicture}
  }
  \caption{The Ising field theory is the continuum scaling limit of
  long-distance spin correlators.  The parameter \(m\) is the mass gap in the
  limiting continuum units, and \(m=0\) gives the critical conformal theory.}
  \label{fig:volume-ii-ising-scaling-limit}
\end{figure}

The operator dictionary at the critical fixed point begins with
\[
  \sigma(x)\longleftrightarrow \phi(x),
  \qquad
  \varepsilon(x)\longleftrightarrow \phi^2(x).
\]
Here \(\sigma\) is the spin field, and \(\varepsilon\) is the energy operator,
the scalar operator coupled to the temperature deformation.  The normalization
of either field is conventional; their scaling dimensions are invariant data.
The relation between the energy dimension and the exponent \(\nu\) is
\[
  \Delta_\varepsilon
  =
  D-\frac1\nu .
\]

\section{A Lattice Scalar Path Integral}

The Ising variables may be enlarged from two values to a real local variable
without changing the long-distance fixed point, provided the relevant
\(\mathbb Z_2\)-even coupling is tuned.  Let \(P:\mathbb R\to\mathbb R\) be an
even single-site potential with sufficient growth at infinity.  The generalized
Ising partition function is
\[
  \widetilde Z_\Lambda
  =
  \int\prod_{x\in\Lambda}\dd s_x\,
  \exp\!\left[-P(s_x)\right]\,
  \exp\!\left[\beta\sum_{\langle x y\rangle}s_xs_y\right].
\]
Equivalently,
\[
  \widetilde Z_\Lambda
  =
  \int\prod_{x\in\Lambda}\dd s_x\,\ee^{-S_\Lambda(s)}
\]
with
\[
  S_\Lambda(s)
  =
  \frac{\beta}{2}\sum_{\langle x y\rangle}(s_x-s_y)^2
  +
  \sum_{x\in\Lambda}\bigl(P(s_x)-\beta D\,s_x^2\bigr),
\]
where each positive lattice direction is counted once.  The identity follows
from
\[
  \sum_{\langle x y\rangle}(s_x-s_y)^2
  =
  2D\sum_xs_x^2-2\sum_{\langle x y\rangle}s_xs_y
\]
on a periodic hypercubic lattice.

For an infinite hypercubic lattice with unit spacing, introduce the Brillouin
zone
\[
  B=[-\pi,\pi]^D
\]
and write
\[
  s_x
  =
  \int_B\frac{\dd^Dk}{(2\pi)^D}\,
  \widetilde\phi(k)\,\ee^{\ii k\cdot x},
  \qquad
  \widetilde\phi(k+2\pi n)=\widetilde\phi(k),
  \quad n\in\mathbb Z^D .
\]
The cutoff \(|k^\mu|\leq\pi\) is a finite ultraviolet cutoff.  The nearest
neighbor kinetic term becomes
\[
  \sum_{\langle x y\rangle}(s_x-s_y)^2
  =
  \int_B\frac{\dd^Dk}{(2\pi)^D}\,
  \widetilde\phi(-k)\widetilde\phi(k)
  \sum_{\mu=1}^D\left|\ee^{\ii k_\mu}-1\right|^2 .
\]
For small momentum,
\[
  \sum_{\mu=1}^D\left|\ee^{\ii k_\mu}-1\right|^2
  =
  k^2-\frac1{12}\sum_{\mu=1}^D k_\mu^4+O(k^6).
\]
Thus the long-distance action has the form
\[
  S[\phi]
  =
  \int \dd^Dx
  \left[
    \frac{\beta}{2}\sum_{\mu=1}^D(\partial_\mu\phi)^2
    +
    V_\beta(\phi)
    -
    \frac{\beta}{24}\sum_{\mu=1}^D(\partial_\mu^2\phi)^2
    +\cdots
  \right],
\]
with \(V_\beta(\phi)=P(\phi)-\beta D\phi^2\) plus local terms generated by the
field normalization convention.  The higher-derivative terms reflect the
lattice cutoff and the breaking of continuous rotations to the hypercubic
lattice symmetry.

\begin{figure}[H]
  \centering
  \resizebox{0.98\textwidth}{!}{%
  \begin{tikzpicture}[x=1cm,y=1cm,every node/.style={font=\scriptsize}]
    \begin{scope}
      \node[font=\small] at (0,2.85) {single-site measure};
      \draw[-{Latex[length=2mm]}, thick] (-1.4,0)--(1.45,0) node[below] {\(s\)};
      \draw[-{Latex[length=2mm]}, thick] (0,-0.25)--(0,2.15) node[left] {\(P(s)\)};
      \draw[blue!70!black, very thick]
        plot[smooth] coordinates {(-1.3,1.9) (-0.75,0.35) (0,1.15) (0.75,0.35) (1.3,1.9)};
      \node[blue!70!black] at (0,-0.55) {\(\mathbb Z_2\)-even};
    \end{scope}

    \begin{scope}[shift={(4.2,0)}]
      \node[font=\small] at (0,2.85) {cutoff momentum space};
      \draw[thick] (-1.3,-1.3) rectangle (1.3,1.3);
      \draw[-{Latex[length=2mm]}, gray!70] (-1.65,0)--(1.8,0) node[below] {\(k^1\)};
      \draw[-{Latex[length=2mm]}, gray!70] (0,-1.65)--(0,1.8) node[left] {\(k^2\)};
      \node at (-1.3,-1.55) {\(-\pi\)};
      \node at (1.3,-1.55) {\(\pi\)};
      \node[align=center, blue!70!black] at (0,-2.0)
        {Brillouin zone \(B\)};
    \end{scope}

    \begin{scope}[shift={(8.8,0)}]
      \node[font=\small] at (0,2.85) {derivative expansion};
      \node[draw=black, rounded corners, align=center, inner sep=7pt] at (0,0.85)
        {\(\displaystyle
        \frac{\beta}{2}(\partial\phi)^2+V_\beta(\phi)
        \)\\[2pt]
        \(\displaystyle
        -\,\frac{\beta}{24}\sum_\mu(\partial_\mu^2\phi)^2+\cdots
        \)};
      \node[align=center, purple!75!black] at (0,-1.05)
        {irrelevant lattice\\anisotropy};
    \end{scope}

    \draw[-{Latex[length=2mm]}, thick] (1.85,0.75)--(3.05,0.75);
    \draw[-{Latex[length=2mm]}, thick] (5.9,0.75)--(7.25,0.75);
  \end{tikzpicture}
  }
  \caption{The generalized Ising model may be rewritten as a scalar lattice
  path integral.  The single-site measure supplies the potential, the nearest
  neighbor interaction supplies the kinetic term, and the Brillouin zone gives
  a finite ultraviolet cutoff.}
  \label{fig:volume-ii-ising-lattice-scalar}
\end{figure}

\section{Universality as Fixed-Point Attraction}

The lattice scalar action just constructed and the continuum quartic scalar
action differ in their short-distance coordinates.  The lattice action has a
finite cutoff, lattice-anisotropic derivative terms, and many allowed local
couplings such as \(\phi^6\) and higher powers or derivatives.  These
differences define distinct points in the same large theory space.  In
\(D=2,3\), after the relevant \(\mathbb Z_2\)-even coupling is tuned to
criticality, their long-distance trajectories approach the same Ising fixed
point.

Near the fixed point, write a general local deformation as
\[
  \Delta S
  =
  \int\dd^Dx\,
  \left(
    h\,\sigma(x)
    +
    \delta m^2\,\varepsilon(x)
    +
    \sum_I\delta g_I\,O_I(x)
  \right),
  \qquad
  \Delta_I=\dim O_I .
\]
The coupling \(\delta g_I\) has mass dimension \(D-\Delta_I\).  At a physical
mass scale \(\mu\), the dimensionless strength of this perturbation is
\[
  \delta g_I\,\mu^{\Delta_I-D}.
\]
If \(\Delta_I>D\), this strength tends to zero as \(\mu\to0\).  These are the
irrelevant perturbations.  They account for microscopic differences such as
the choice of square or hexagonal lattice, the details of \(P(s)\), and
higher-derivative cutoff terms.

The two relevant scalar directions of the Ising fixed point are the spin field
\(\sigma\) and the energy field \(\varepsilon\):
\[
  \Delta_\sigma<D,
  \qquad
  \Delta_\varepsilon=D-\frac1\nu<D .
\]
The magnetic field \(h\) couples to \(\sigma\).  The temperature deviation
couples to \(\varepsilon\):
\[
  \Delta S_T
  =
  \int\dd^Dx\,\delta m^2\,\varepsilon(x),
  \qquad
  \delta m^2\propto T-T_c .
\]
Since \(\delta m^2\) has dimension \(D-\Delta_\varepsilon\), it generates the
correlation length
\[
  \xi(T)
  \propto
  |\delta m^2|^{-1/(D-\Delta_\varepsilon)}
  \propto
  |T-T_c|^{-\nu}.
\]
This relation is the RG origin of the exponent \(\nu\).

\begin{figure}[H]
  \centering
  \resizebox{0.96\textwidth}{!}{%
  \begin{tikzpicture}[x=1cm,y=1cm,every node/.style={font=\scriptsize}]
    \node[font=\small] at (0,3.6) {theory space near the Ising fixed point};
    \fill[purple!80!black] (0,0) circle (3pt);

    \node[draw=black, rounded corners, align=center, fill=blue!7, inner sep=5pt]
      (phi) at (-4.2,2.2)
      {continuum scalar\\quartic actions};
    \node[draw=black, rounded corners, align=center, fill=green!7, inner sep=5pt]
      (lat) at (4.2,2.2)
      {Ising lattice\\actions};
    \node[draw=black, rounded corners, align=center, fill=gray!8, inner sep=5pt]
      (irr) at (0,2.75)
      {irrelevant local\\differences};

    \draw[-{Latex[length=2mm]}, very thick, blue!75!black]
      (phi.south east)..controls (-2.5,1.2) and (-1.3,0.55)..(-0.15,0.1);
    \draw[-{Latex[length=2mm]}, very thick, green!50!black]
      (lat.south west)..controls (2.5,1.2) and (1.3,0.55)..(0.15,0.1);
    \foreach \x in {-1.5,-0.7,0.7,1.5} {
      \draw[-{Latex[length=2mm]}, gray!70, thick]
        (\x,2.25)..controls (0.6*\x,1.55) and (0.25*\x,0.8)..(0.05*\x,0.18);
    }

    \draw[-{Latex[length=2mm]}, red!75!black, very thick]
      (0,0)--(0,-2.0);
    \node[red!75!black, align=center] at (1.85,-1.15)
      {temperature or mass\\deformation};
    \node[blue!70!black, align=center] at (-2.55,-1.35)
      {critical tuning\\sets this coordinate to zero};
    \draw[-{Latex[length=2mm]}, blue!70!black] (-1.8,-1.1)--(-0.2,-0.25);
    \node[purple!80!black, align=center, fill=white, inner sep=2pt,
      anchor=east]
      at (-0.25,0.72)
      {Ising CFT\\fixed point};
  \end{tikzpicture}
  }
  \caption{Universality is fixed-point attraction in theory space.  Lattice
  spin models and scalar field theories can have different ultraviolet
  coordinates while their critical trajectories approach the same Ising fixed
  point.  Relevant deformations move the theory away from criticality and
  produce massive scaling limits.}
  \label{fig:volume-ii-ising-universality-rg}
\end{figure}

Thus the statistical Ising model supplies a precise bridge from microscopic
lattice ensembles to continuum Euclidean QFT.  The bridge is made by
correlators: after tuning the relevant directions and rescaling fields, the
large-distance spin correlators become correlators of local fields at, or near,
the Ising fixed point.
